% =========================================================================
% SECCIÓN 1: INTRODUCCIÓN
% =========================================================================

\section{Introducción}
\label{sec:introduccion}

\subsection{Contexto y Motivación de la Revisión}
\noindent [Descripción del contexto tecnológico global y regional, tendencias en Small Language Models (SLMs), arquitecturas RAG y la brecha digital en lenguas originarias como el Quechua Collao.]

\subsection{Planteamiento del Problema de Revisión}
\noindent [Formulación del problema o vacío de conocimiento en la literatura que justifica la presente revisión sistemática, vinculando las limitaciones de cómputo en hardware accesible con la falta de recursos lingüísticos computacionales.]

\subsection{Objetivos de la Revisión Sistemática}

\subsubsection{Objetivo general}
\noindent [Definición formal del objetivo general de la revisión sistemática de la literatura.]

\subsubsection{Objetivos específicos}
\begin{enumerate}[(a)]
	\item [Objetivo específico 1: Caracterizar arquitecturas de modelos compactos de lenguaje aptas para inferencia local offline.]
	\item [Objetivo específico 2: Analizar mecanismos de recuperación aumentada adaptados a lenguas de bajos recursos y morfología aglutinante.]
	\item [Objetivo específico 3: Evaluar estrategias de adaptación eficiente en parámetros (PEFT/QLoRA) para dominios educativos interculturales.]
	\item [Objetivo específico 4: Sintetizar las brechas tecnológicas y metodológicas del estado del arte para formular una hoja de ruta de investigación.]
\end{enumerate}

\subsection{Alcance y Delimitación de la Revisión}
\noindent [Definición del alcance temático, marco temporal cubierto (últimos 5 años) y delimitación lingüística centrada en Quechua Collao y lenguas andinas afines.]

\subsection{Estructura del Artículo}
\noindent [Descripción concisa de la organización de las secciones restantes del artículo científico.]
